\section{Lebesgue 积分先积简单函数再逼近}
\label{sec:12-Real-Analysis-Support/05-Measure-Integration-and-Conditional-Expectation/03}

试着用 Riemann 积分算这个函数在 \([0,1]\) 上的面积：

\[
f(x)=\ind_{\Q\cap[0,1]}(x)=
\begin{cases}
1,&x\in\Q,\\
0,&x\notin\Q.
\end{cases}
\]

按 Riemann 的办法，你把 \([0,1]\) 切成很多小区间，在每个区间上找函数值的上下界，用矩形近似。问题是：随便多小的区间，里面都同时住着有理数和无理数——有理数在其中稠密，无理数也是。所以每个子区间上的上确界永远是 \(1\)，下确界永远是 \(0\)。上和恒为 \(1\)，下和恒为 \(0\)，再怎么切也没有收敛的趋势。Riemann 积分告诉你：这个函数不可积。

但直觉上你可能会觉得它``应该''积分为零——毕竟无理数的数量``远远多于''有理数，值为 \(1\) 的点只占了实数线上微不足道的一丁点。Riemann 积分没有``测度''这个概念来量化``一丁点''，所以无法把直觉兑现成积分数值。

\subsection{换一个组织方式：先按高度分层}

Riemann 积分把定义域切成竖条（小区间），每条的宽度是 \(\Delta x\)，高度取函数值，用矩形面积逼近图像下的区域。

Lebesgue 的想法是横着来。函数在 \([0,1]\) 上的取值只有 \(0\) 和 \(1\) 两个高度。先看值为 \(1\) 的那一层：哪些 \(x\) 映射到了 \(1\)？有理数集合。这层有多大？它虽然是无穷个点，但在实数线上的 Lebesgue 测度为零。值为 \(0\) 的那一层占了多少？几乎全部：无理数集合的测度是 \(1\)。

积分就是一种``高度乘以该高度占据的测度''的累加。高度 \(1\) 占据测度 \(0\)，贡献 \(1\times0=0\)。高度 \(0\) 占据测度 \(1\)，贡献 \(0\times1=0\)。总和为零。

Dirichlet 函数在 Riemann 手下不可积，Lebesgue 手下积分就是零。Lebesgue 没有改变被积函数——它改变的是组织求和的方式：不再切定义域，而是按取值分层，测量每层原像的测度，再求和。这个方向对概率论特别自然，因为随机变量的取值层和事件正是通过原像连接的：\(\Pr(X\in B)=P(X^{-1}(B))\)。

\subsection{从能直接算的函数出发：先学会走}

对可测集合 \(A\)，指示函数的积分就是该集合的测度：

\[
\int_\Omega \ind_A\,d\mu=\mu(A).
\]

这是整个 Lebesgue 积分的原点——没有比这更简单的函数了，只取零和一，积分只依赖于支撑集的测度。

往前走一步。若非负简单函数只取有限多个值，可以写成各层指示函数的加权和：

\[
s=\sum_{k=1}^m a_k\ind_{A_k},
\qquad a_k\ge0,
\]

其中 \(A_k\) 两两不交——每个点只属于一个层。积分就是各层高度乘以该层测度：

\[
\int_\Omega s\,d\mu
=\sum_{k=1}^m a_k\mu(A_k).
\]

一个简单函数可以用不同的分层方式表达：比如把某层一劈为二，两个子层的值写成同样的高度。但只要回溯到所有层的最细共同分割，两种表达在该分割上的求和结果必然一致。测度的可数可加性保证了积分值不依赖你选择的分层表述。

\subsection{通到一般函数的桥：用简单函数从下面顶上去}

你已经会给简单函数积分了。对于一般的非负可测函数 \(f\)，想法是：找到所有不超过 \(f\) 的简单函数，它们的积分构成一个集合。\(f\) 的积分就是这个集合的上确界：

\[
\int_\Omega f\,d\mu
=\sup\set{\int_\Omega s\,d\mu:
0\le s\le f,\ s\text{ 为简单函数}}.
\]

这个积分允许等于 \(\infty\)。不是出 bug——如果函数在某些区域贡献了无限的面积，积分就应该反映这个事实。

构造性地看这个过程。对非负可测函数，你总可以造出一列递增的简单函数 \(s_n\uparrow f\)。例如把值域从 \(0\) 到 \(n\) 切成 \(n\cdot2^n\) 个等宽小层，低于 \(n\) 的取值落进哪一层就取哪一层的下限，高于 \(n\) 的截断为 \(n\)。随着 \(n\) 增大，\(s_n\) 从下面越来越贴近 \(f\)。你不需要对复杂的 \(f\) 直接写积分——你在用一列你会算积分的简单函数，从下往上逼近它。上确界定积分等价于取这个递增序列的积分极限。

这是 Lebesgue 积分的构造入口：把``简单函数的积分''作为积木块，把``上确界/递增极限''作为搭桥手段。所有非负可测函数从此都有了积分定义。

\subsection{正负分家：别让无穷大互相抵消}

一般实值函数不一定非负。把它掰成正负两半：

\[
f=f^+-f^-,
\qquad
f^+=\max(f,0),
\quad
f^-=\max(-f,0).
\]

两部分都是非负的，各自的积分已经有了定义。如果至少一个有限，就可以相减：

\[
\int f\,d\mu
=\int f^+\,d\mu-\int f^-\,d\mu.
\]

若正部分和负部分的积分都是 \(\infty\)，就遇到了 \(\infty-\infty\)——一个没有定义的形式。不能因为函数在正负两侧对称就宣布积分是零，因为无穷大是无法靠对称性抵消的量。

通常说``函数可积''（integrable）是指

\[
\int\abs{f}\,d\mu<\infty,
\]

这等价于正负部分积分都有限。Cauchy 分布的均值不存在，正是这个原因：两侧尾部太厚，正半轴和负半轴分别贡献了无穷大的面积，\(\infty-\infty\) 不是零，是没有定义。

\subsection{期望就是概率测度下的 Lebesgue 积分}

现在重新看随机变量的期望。\(\E[X]\) 的定义是

\[
\E[X]=\int_\Omega X\,dP,
\]

没有什么附加条件——这就是概率测度 \(P\) 下的 Lebesgue 积分。概率空间上的可测函数就是随机变量，可积性条件 \(\E[\abs{X}]<\infty\) 就是可积定义的重述。

你若只知道分布 \(P_X\) 而不直接接触底层的样本空间 \(\Omega\)，换元公式用推前测度把积分从 \(\Omega\) 搬运到 \(\R\)：

\[
\int_\Omega g(X(\omega))\,dP(\omega)
=\int_\R g(x)\,dP_X(x).
\]

右侧的形式因参考测度不同而显现为不同面貌。离散随机变量用计数测度作为参考：

\[
\sum_x g(x)P_X(\set{x});
\]

连续型具有密度时用 Lebesgue 测度作为参考：

\[
\int_\R g(x)p(x)\,dx.
\]

初等概率论里这是两个独立定义——一个叫求和，一个叫带密度的积分。测度论把它们统一了：两者都是同一个式子的不同写法，差别只在于你选用哪个参考测度来计量``测度分配''。概率论里的离散和连续不再是一道非此即彼的分岔口，而是一个统一框架里的两种测度选择。

\subsection{Lebesgue 积分对极限友好，因为它的构造就是为了极限}

Riemann 积分的困难常常出在极限上：一列 Riemann 可积函数的逐点极限，可能不再 Riemann 可积（Dirichlet 函数就是一列简单函数——从 \(0\) 到 \(n\) 逐层逼近——的极限，但极限本身不可 Riemann 积）。Lebesgue 积分的构造方式——从简单函数通过 sup/极限往上搭——天然让它对极限运算更从容。

三个定理管住极限与积分的交换。

非负可测函数 \(f_n\uparrow f\) 时，单调收敛定理保证

\[
\int f_n\,d\mu\uparrow\int f\,d\mu.
\]

不需要附加条件——非负性和单调性就足够了。Fatou 引理在失去单调保证时仍给了 liminf 下界：

\[
\int\liminf_{n\to\infty} f_n\,d\mu
\le\liminf_{n\to\infty}\int f_n\,d\mu.
\]

控制收敛定理则在最强的假设下给出最干净的结果：若 \(f_n\to f\) 几乎处处，且存在可积函数 \(g\) 使 \(\abs{f_n}\le g\)，则

\[
\lim_{n\to\infty}\int f_n\,d\mu
=\int f\,d\mu.
\]

三条定理从弱到强覆盖了极限取进积分的大多数场景。它们不是 Lebesgue 积分造完之后追加的补救，而是 Lebesgue 积分的构造本身——简单函数逼近加测度的可数可加性——推导出来的直接后果。

\subsection{双重积分的交换同样需要一份许可}

乘积测度空间上有两重积分摆在面前时，你几乎总会想交换积分次序。交换不是想换就换。

若被积函数非负可测，Tonelli 定理无条件放行：

\[
\int\!\left(\int f(x,y)\,d\nu(y)\right)d\mu(x)
=
\int\!\left(\int f(x,y)\,d\mu(x)\right)d\nu(y),
\]

两边可以同时是 \(\infty\)。非负性是通行证——只要函数不取负值，交换次序不会引入条件收敛造成的次序依赖。

若函数有正有负，Fubini 定理要求绝对可积：

\[
\int\!\int\abs{f(x,y)}\,d\mu(x)d\nu(y)<\infty.
\]

没有这条线，不同的积分次序可以走向不同的有限极限值，每个对应不同的条件收敛抵消路径。这不是定理太苛刻——是条件收敛本就不是一个秩序无关的量。概率论中交换期望、交换无穷求和与积分、或交换多个坐标的积分时，应该先说清楚自己是在用 Tonelli（非负性）还是 Fubini（绝对可积），而不是笼统一句``积分可以换序''。
